\input{../tex-inputs/latex-leseansicht-vorspann}

\section[Arthur Schnitzler an Rainer Maria Rilke, {[}zwischen 20. 5. und 20. 6. 1901?{]}]{L04346 Arthur Schnitzler an Rainer Maria Rilke, {[}zwischen 20. 5. und 20. 6. 1901?{]}}
\nopagebreak\mylabel{L04346v}
\rehead{ }\normalsize\beginnumbering\briefempfaengerindex{Rilke, Rainer Maria@\textsc{Rilke, Rainer Maria}!zzzSchnitzler, Arthur@\emph{von Arthur Schnitzler}!1901-06-202@{{[}zwischen 20. 5. und 20. 6. 1901?{]}}|(be}
\toendnotes[C]{\smallbreak\pagebreak[2]}
\correspDesc{Versand  durch Arthur Schnitzler im Zeitraum [zwischen 20. 5. und 20. 6. 1901?] in Wien
\newline{}Erhalt  durch Rainer Maria Rilke im Zeitraum [zwischen 21. 5. und 21. 6. 1901?] in Dresden}\toendnotes[C]{\smallbreak}
\Standort{DLA, A:Schnitzler, HS.2022.81.}
\physDesc{Briefkarte, 93 Zeichen
\newline{}Handschrift: schwarze Tinte, deutsche Kurrent}
\buchAbdrucke{\weitereDrucke{\emph{Rainer Maria Rilke und Arthur Schnitzler. Ihr Briefwechsel.}. Mit Anmerkungen herausgegeben von Heinrich Schnitzler In: \emph{Wort und Wahrheit}, Jg. 13, Nr. 4, April 1958, S. 282–298, hier S. 292.} }\toendnotes[C]{\smallbreak}
\pstart
           \noindent{}{\pb}Herzlichen \label{K_L04346-1v}\edtext{Glückwunſch}{\lemma{\textnormal{\emph{Glückwunsch}}}\Cendnote{\textnormal{Die Karte
                  ist undatiert. Unter der Annahme, dass sie im Zeitraum der aktiven Korrespondenz zwischen Rilke\pwindex{Rilke, Rainer Maria 4.\,12.\,1875 Prag – 29.\,12.\,1926 Montreux@\textsc{Rilke, Rainer Maria} (4.\,12.\,1875 Prag – 29.\,12.\,1926 Montreux)|pwk} und Schnitzler
                  (1896–1902) versandt wurde, bietet sich vor allem die Verehelichung Rilkes\pwindex{Rilke, Rainer Maria 4.\,12.\,1875 Prag – 29.\,12.\,1926 Montreux@\textsc{Rilke, Rainer Maria} (4.\,12.\,1875 Prag – 29.\,12.\,1926 Montreux)|pwk} mit Clara Westhoff\pwindex{Westhoff, Clara 21.\,11.\,1878 Bremen – 9.\,3.\,1954 Fischerhude@\textsc{Westhoff, Clara} (21.\,11.\,1878 Bremen – 9.\,3.\,1954 Fischerhude)|pwk} am 28. 4. 1901\eventindex{Bremen@\textbf{Bremen}!Hochzeit von Clara Westhoff und Rainer Maria Rilke, 28.4.1901@Hochzeit von Clara Westhoff und Rainer Maria Rilke, 28.4.1901|pwk}
                  an, auf die Rilke\pwindex{Rilke, Rainer Maria 4.\,12.\,1875 Prag – 29.\,12.\,1926 Montreux@\textsc{Rilke, Rainer Maria} (4.\,12.\,1875 Prag – 29.\,12.\,1926 Montreux)|pwk} in seinem Brief vom XXXX Auszeichnungsfehler: Dokument L04538 nicht gefunden anspielt. Schnitzler
                  dürfte diese Gratulationskarte einem Widmungsexemplar von \emph{Lieutenant Gustl}\pwindex{Schnitzler, Arthur 15. 5. 1862 Wien – 21. 10. 1931 ebd.@\textsc{Schnitzler, Arthur} (15. 5. 1862 Wien – 21. 10. 1931 ebd.)!Lieutenant Gustl. Novelle@\strich\emph{Lieutenant Gustl. Novelle}|pwk} beigelegt haben. Diese standen ihm
                  wiederum circa ab dem 20. 5. 1901 zur Verfügung, so dass die vorliegende Karte frühestens um diese Tage und spätetens
                  kurz vor dem Erhalt des Antwortschreibens vom XXXX Auszeichnungsfehler: Dokument L04537 nicht gefunden anzusetzen ist.}}}\label{K_L04346-1}, verehrter Herr Rilke,{ }ſagt Ihnen\pend
           
\pstart
           Ihr aufrichtig ergebner{\\[\baselineskip]}\spacefill\mbox{Arthur Schnitzler}\pend
           \leftskip=0em{}\selectlanguage{ngerman}\endnumbering\briefempfaengerindex{Rilke, Rainer Maria@\textsc{Rilke, Rainer Maria}!zzzSchnitzler, Arthur@\emph{von Arthur Schnitzler}!1901-05-202@{{[}zwischen 20. 5. und 20. 6. 1901?{]}}|)be}\mylabel{L04346h}
\begin{anhang}
\end{anhang}\newcommand{\dateiname}{L04346}\newcommand{\titel}{Arthur Schnitzler an Rainer Maria Rilke, [zwischen 20. 5. und 20. 6. 1901?]}\newcommand{\editorInnen}{Herausgegeben von Jahnke, SelmaMüller, Martin Anton}\input{../tex-inputs/latex-leseansicht-abspann}
